\chapter{Construction des tables offres, metiers et competences}

\section{Modèle relationnel}

Trois tables construites dans \texttt{src/etl/tables.py} à partir du
DataFrame nettoyé (chapitre~6), écrites en Parquet sur HDFS
(\texttt{/projet\_ia/processed/}) puis exposées comme tables Hive
externes (\texttt{processed.offres}, \texttt{processed.metiers},
\texttt{processed.competences}) via \texttt{src/etl/hive\_io.py}.

\begin{itemize}[leftmargin=1.5em]
  \item \textbf{metiers} — référentiel : une ligne par
    \texttt{job\_title} distinct (\texttt{job\_title},
    \texttt{job\_category}, \texttt{metier\_id}).
  \item \textbf{offres} — grain = une ligne par offre
    (\texttt{job\_id}, \texttt{metier\_id}, salaires, localisation,
    expérience, \dots), jointe au référentiel \texttt{metiers}.
  \item \textbf{competences} — une ligne par couple
    (\texttt{job\_id}, compétence), obtenue par
    \texttt{split(required\_skills, "\textbackslash{}|")} puis
    \texttt{explode()}.
\end{itemize}

\begin{constat}[Correction : \texttt{metier\_id} déterministe]
La première version (issue du notebook d'exploration) générait
\texttt{metier\_id} avec
\texttt{monotonically\_increasing\_id()}. Cette fonction Spark
\textbf{n'est pas stable} entre deux exécutions ni entre partitions
(elle encode le numéro de partition dans les bits hauts de l'ID) : deux
exécutions du pipeline sur le même dataset peuvent produire des
\texttt{metier\_id} différents pour le même métier, ce qui casserait
silencieusement toute jointure historique ou tout ré-exécution
partielle du pipeline.

Corrigé en \texttt{src/etl/tables.py} :
\texttt{metier\_id = substring(sha2(job\_title, 256), 1, 16)} — un
hash déterministe de \texttt{job\_title}, identique à chaque exécution
et à travers les partitions. Vérifié par le test
\texttt{tests/unit/test\_tables.py::test\_build\_metiers\_is\_deterministic\_and\_deduplicated}
qui exécute \texttt{build\_metiers()} deux fois sur les mêmes données
et compare les ID obtenus.
\end{constat}

\section{Chiffres réels (run du 16/09/2026)}

\begin{center}
\begin{tabular}{ll}
\toprule
Table & Nombre de lignes \\
\midrule
\texttt{processed.offres} & 1500 \\
\texttt{processed.metiers} & 25 (métiers distincts) \\
\texttt{processed.competences} & 9548 (liaisons job\_id $\times$ compétence) \\
\midrule
Doublons \texttt{job\_id} (avant construction) & 0 \\
\bottomrule
\end{tabular}
\end{center}

Ces chiffres sont cohérents entre eux : 1500 offres pour 25 métiers
distincts (60 offres en moyenne par métier), et 9548/1500 $\approx$
6,4 compétences en moyenne par offre.
